\subsubsection{B 类：变量、约束与策略必须逐项对应}

B 类先写决策对象、目标统计定义和约束来源。典型信息链是“令……表示是否选择……”
“由产能/路径/时序要求得到约束……”“为比较冲突目标，构造……并保留 Pareto
方案”“在相同函数评估预算下，改进方法相对基线……”。每个二元变量应能翻译回
一个实际动作，每个约束都能回到题意，每个输出都能转化为方案表。

不能只给元启发式流程图和最优值。论文需报告编码、初始化、约束修复、参数、
终止条件、随机种子或多次运行统计；能使用精确规划时，还应报告界或间隙。图表
以决策结构、算法流程、收敛/Pareto 图、方案对照和敏感性图为主。

B078、B108、B125、B175、B007、B026、B050、B160、B030、B035、B086、B226、B311、
B477、B159、B195、B196、B060 与 B157 共十九篇 B 题均已通过人工检查。跨论文反复出现的
写作骨架是“信息集—状态—动作—转移—目标—求解—检验”：
已知信息的单人层先给最短路、动态规划或整数规划基线；未知信息层说明概率来源、
情景生成器、缺货风险和滚动再决策，并把知道完整未来的离线全知上界与只使用当时
信息的在线非预见策略分开；多人层先判定共同收益、一般和、零和或拥塞结构，再定义
联合状态和支付，依次检查最佳反应、支配关系、纯均衡、内点混合解与边界解；若是
合作多目标，则明确加权、词典序、约束法或 Pareto 选择。终点
残值必须按题意进入目标，整数库存和路线候选须覆盖完整可行域；候选集上的混合均衡
不宜写成完整路径空间的确定最优路线，逐日局部均衡也不能替代跨日状态耦合的全局策略。

相应结果段应把三类证据分开：确定性层报告目标值、库存余量和约束满足；随机层
报告独立重复数、种子、均值、标准差、分位数与失败率，并使类别计数、筛选后分母
和比例在舍入容差内互相还原；多人层报告支付如何映射到题面收益、候选集覆盖和各
玩家可行性。灵敏度分析还须声明每个扰动点是固定原策略
复算还是重新优化，并同时记录目标变化、策略切换和活跃约束。

B007 给出了另一种常见的 B 类数据--优化混合链：先按受控条件整理实验数据，利用
Pearson 相关与二项式拟合辨认单因素关系，再用 Q 型聚类归并催化剂组合；随后在同一
数据和同一指标下比较多元回归与偏最小二乘，针对原模型未直接描述产物收率这一
缺口建立收率回归，并增加入温度交互项，最后把回归关系送入 LINGO 优化，以均匀
设计补充有限预算下的实验点。这里的优秀之处不是模型多，而是每次增加复杂度前都
先指出上一环节没有回答什么，下一环节又有明确输入和输出。

其结果段也采用容易核读的成对结构：先报告无温度限制下的最优值，再写温度不高于
$350\,^{\circ}\mathrm{C}$ 时的受限最优值，并说明为什么工程上需要第二个方案。实验设计段
先写“现有两种方法可以选择”，逐一说明适用性和舍弃理由，并落到均匀设计及安全
修正。读者看到的是一次真实取舍，而不是事后把唯一方案包装成必然选择。

B026 把这条链继续向“证据缺口驱动实验”推进。作者先用三次样条统一温度水平，
再从散点趋势和转化率上界说明 Logistic 的平台形态；选择性和时间曲线则分别保留
二次关系与分段下降后的稳定段。双因素方差分析先回答温度、催化剂组合是否有影响，
多重比较再定位具体水平，随后把催化剂配方改写为连续变量，分别建立两种装料方式的
二次回归并在全温区、温度不高于 $350\,^{\circ}\mathrm{C}$ 两种统计定义下求解。这里的
人类写法在于每一步都给出“前一张表还没有回答什么”：缺失组合补测、未观测区间补测、
保持单一变量的对照和最终五次新增实验，分别对应不同的裁决问题；并非笼统声称“增加
实验以验证模型”。

该篇还留下可复核用的结果句式：先以 F 值和 $p$ 值判断因素是否有影响，再以均值差
回答哪个水平更优；装料方式 I/II 的拟合优度分别为 $90.16\%/90.78\%$，四种优化
情景保持同一目标统计定义，结果表中的温度为 $400/349/400/349\,^{\circ}\mathrm{C}$。
这些数字只作为论文报告值的页码证据，写作时仍需说明比较条件和结果用途。

B050 又补出两种很像真实建模过程的推进方式。其一，候选模型并不是按拟合优度机械
排序：四次多项式在五个样本点上得到 $R^2=1$，作者随即指出这只是把有限样本完全
插值，无法据此判断整体规律，因而改用三次多项式。这样一段同时留下候选结果、样本
限制和最终裁决，比“综合考虑后选择三次模型”容易读者可以相信取舍确实发生过。

其二，递进回归每次都保留前一轮指标和尚未解决的问题。基础结构的决定系数/调整后
决定系数为 $0.795433/0.685442$，加入非线性变换后变为 $0.828916/0.696456$，再加入
交互项后达到 $0.934994/0.802786$。这些数字不是并排陈列，而是分别承担“现有模型还
缺什么”和“新增结构是否值得”的判断。随后，作者先写多项式回归送入规划后结果不理想，
才改用 BP 网络；方法切换由具体失败触发，不用“为提高精度”一笔带过。最后五次新增
实验先分为三次峰值探索和两次预测确认，再列实验点及用途，读者因此知道每个点是在
扩展搜索范围还是核对模型读数。

B160 把“控制变量法”和“回归效果较好”这两类常被写空的段落补成了可核读结构。
控制变量比较并不只报因素名称，而是先列其余保持条件和唯一改变量，点名可比组，
随后在同一温度下分别读取乙醇转化率与 C4 选择性，最后只归结当前组的局部结论。
因此，读者既能核对组间是否可比，也不因单个响应改善就误把配方写成总体占优。

回归段则保留“初始模型--诊断图--多重共线性--逐步回归--修正模型”的真实轨迹。
转化率模型依次报告 $R=0.91$、$R^2=0.832$、$RMSE=0.0417$、$F=39.74$ 与
$p=6.17\times10^{-27}$；选择性模型相应为 $0.89$、$0.80$、$0.0453$、$36.03$ 与
$1.14\times10^{-24}$。作者把前两项称作相对效果，把 RMSE 作为绝对效果，再由 $F/p$
说明整体显著性，三层证据各有任务。问题三又按全温区和不高于 $350\,^{\circ}\mathrm{C}$
重新求解，论文报告收率分别为 $56.3\%$ 与 $29.8\%$；D1--D5 则分别承担曲线延拓、
缺失温区、两类响应水平组合和优化点实测的裁决，不以一句“增加五次实验验证模型”
掩去用途差异。

B030 的写法从“先求全”开始。论文先将三架无人机的相对位置列成三类九种形态，
随后才写“在求解过程中，发现”其中两类可通过旋转重叠，于是把重复推导压成六个
定位分支。这里的简化不是作者预先宣告的愿望，而是完整分类后观察到的几何事实。
算法段同样不写空泛的“进行仿真”：先固定发射机和极径，再列另一发射机、接收机、
六种模型及两个 $[-1^\circ,1^\circ]$ 误差网格，最后写“记录所得全部结果，并绘图输出”。
对象、范围和产物都在同一段内闭合。

该篇还给出两种长文中自然的节制。10500 组结果不硬塞进正文，而是先说明完整结果
保存于 Excel 支撑材料，正文只展示一个代表情形；锥形编队也不再重写整套公式，而是
先以 FY05 为圆心转成半径 $50\,\mathrm{m}$、$50\sqrt{3}\,\mathrm{m}$ 和
$100\,\mathrm{m}$ 的三层圆周，再写“其调整推导与问题一（3）完全一致，对此不再重述”。
结果段先报最大角度和半径调整量，继而用飞行距离解释两者为何量级不同，再以“虽有误差，
但误差较小”保留让步，最后作整体坐标对照。这样的行文比连续使用“有效地”“充分证明”
更像真实作者，因为省略、复用和判断都有可指出的依据。

B035 让另一种方法切换显得更像真实求解。论文先写极坐标正弦定理的四区域分类，
继而直说“这种分类讨论求解方程的做法过于麻烦”，再以定弦定角对应两圆轨迹改写定位。
新方法并非只换一个名称：旋转矩阵与叉积判定圆心所在侧，两圆心连线给出另一交点的
对称关系，原先的分区方程因此变成可画、可检查的几何构造。这里可迁移的是“旧方法--
可见负担--保留不变量--新构造--减少什么计算”的完整转折。

误差论证也没有用一句“经证明可识别”带过。作者先举出一个具体偏差例，随后将候选位置
限制到轨迹圆与极角扇形相交的短弧；随后检查端点，证明第三个角的变化范围不超过
$2^\circ$，最后用分区角度表穷举无交叠区间。调整算法比较时则承认“在局部上可能
贪心算法占优势”，再解释前视搜索为何会修正误差更大的点，使全局目标反转。六轮
方案的 $L_2$ 损失约为 $10^{-23}\,\mathrm{m}^2$，但正文随即说明实际应用只取前五轮、
共 16 架次；额外数值精度并不自动值得通信与编队代价。问题二复用前问前又集中列出
三点共线、四点共圆等退化情形及其修复。这些段落把让步、停止和失败输入输出关系都留在正文中。

B086 进一步说明，几何优化论文的“信息合法性”也可以写得清楚而自然。论文先由
“无人机间信息不共享”排除直接使用全局队形误差，再让每架接收机只比较自身可见的
角度向量与理想角度向量。随后又定义全局误差，但紧接着说明该量只在仿真结束后用于
评价，各无人机并没有据此调整。于是“局部目标”和“全局评价”不是两个容易混淆的
公式，而是分别回答“现场能用什么信息”和“赛后怎样比较方案”。

径向预处理也没有停在 $O'_j/R\leq1.049$ 这一误差界上。作者先据此把题设
$\pm15\%$ 的半径偏差压到约 $\pm5\%$，再用“更重要的是”说明搜索区域缩小后，局部
穷举需要的信号与计算次数也随之减少。两种调整方案使用同一组原始坐标，先比较收敛
速度，再比较稳定定误差，最后估算搜索发射代理：方案一和方案二分别为 $134400$ 与
$25600$，后者约为前者的 $19\%$。论文又特别指出该代理没有实际调整次数的直接含义，
因此最终裁决同时保留了效果、资源和统计定义边界。

锥形编队段把三条几何引理写成了实际流程。每条引理都交代固定点、调整点、控制角或
比例以及目标图形，随后映射到 A1--A15 的具体机群，组成三阶段调整。平面拟合检验值
由 $3.893$ 降至 $0.206$。稳健性段则先承认单组原始坐标不足以检验模型，再给半径和
角度加入高斯微扰，使用十组数据重复计算，并把结论限定在论文所报 $10^{-7}$--
$10^{-5}$ 收敛量级内。这样的承认、补证和限定，比写“模型具有较强鲁棒性”更像
真实作者的论证。

B226 又补出一类很少被算法清单捕捉到的人工写法。作者先说“由问题一知，只要已知
水深和测线与坡面法线投影的夹角，就可以求覆盖宽度”，于是三维问题没有被包装成
一个新的“综合模型”，而是被压成两个待求量。随后用“为了让模型的建立过程更加
直观易懂，我们记……为面 I，……为面 II”给几何对象命名，才进入法向量、叉乘和
线面夹角。这样的段落先读者可以知道在算什么，再让代数承担计算，比直接堆向量式自然。

测线方向也没有只靠规范背书。正文先引用作业规程说明主测线通常平行等深线，随即
写“出于建模的严谨性，我们在下文证明”，再用命题和同长度候选的梯形/矩形覆盖面积
比较完成授权。测线间距则从东边界锚定第一条测线，以 1 m 为步长向西搜索，只保留
$10\%$--$20\%$ 重叠率并取最低可行值。这里的“以……为始边”“恰好能达到……”“从
第一条测线开始”不是装饰性连接词，而是直接对应初始化、步长、可行域和选择规则。

第四问保留了另一种真实取舍。作者直说“受到问题三模型的启发，我们人为地将海域
划分为四个区域”，并紧跟等深线方向这一依据；局部平面拟合后，区域 III 和 IV 的
均方误差过大，才只拆这两个失败区域，已合格区域保持不动。结果段也不写统一的
“方案合理”：区域 V 因坡度大、单线覆盖窄而测线多；区域 II 因沿用区域 I 的延长线
以降低作业难度而出现集中漏测。异常区域、局部原因和数值后果逐项相接，正是应写入
Skill 的人类段落组织。

B311 对“题意存在两种理解”给出了一条更值得保留的写作路径。作者没有先替题面
选定定义，而是分别建立两个重叠率模型，把同一组测线数值代入，再问两种模型会在
CAD 图中留下什么不同的可观察区域。图形裁决后才舍弃模型一、保留模型二。这样的
“歧义--备选--共同例子--判别量--取舍”把真实犹豫留在正文中，比写“经分析采用
定义二”更像人类的求解记录。

跨问推广也只改真正变化的量。作者先比较本问与问题一的几何结构，保留原覆盖关系，
再重新计算当前位置水深和测线方向与坡面的夹角；函数输入输出关系因此只替换这两个输入。
方向判断没有因为坡度小便立刻忽略，而是把它积累到 2 海里尺度，发现两端覆盖宽度
已经相差数百米、可行间距区间不能兼容后，才改用平行等深线。`GB12327-2022` 出现在
定量裁决之后，只作补充印证。这种句序先说明题内证据，再引外部规范，不让标准替代推导。

模型检查与递推结尾同样具体。“首先、其次、再次、最后”分别核对对称性、单调性、
沿等深线方向的不变性和特殊角度下的极限式，并非为了凑四句排比。问题三从西边界
确定首线，按 $10\%$ 重叠递推到第 34 条，再用东侧剩余 $22.1\,\mathrm{m}$ 和
$3713.7>3704$ 说明第 35 条应舍弃；问题四又以等深线间距最大的最平缓位置作为最不利点，
由该处完全覆盖推出其余位置覆盖。结构预言、边界余量和最不利位置各自承担一个结论，
所以段落虽使用常见连接词，仍有清楚的人类判断痕迹。

B477 把这种真实判断推进到“定义失效--换法--局部返工”三个节点。旧重叠率并非
被一句“存在局限”带过：作者固定相同测线间距和重叠长度，画出直观重叠更大的情形
反而得到更小比例，写明继续使用会使指标失去意义，随后只把间距改成坡面方向上的
投影量。二维关系推广到三维以前，又先证明航向固定时局部坡度唯一，随后将当前位置
深度和等效坡度送回原式。段落的自然感来自反例只触发必要修改，证明也确实授权下一步。

真实海域部分没有用“为提高精度，采用随机森林”替代选择过程。论文先将纯离散加密
的负担拆成精度、局部失效、误差累积以及数据量与耗时四项，再要求学习模块输出任意点
的深度、梯度、坡度和覆盖宽度，因而读者知道新方法交付给哪一段路线算法。第一次搜索
出现起点依赖、局部漏测和可用性差时，正文也保留了这个失败状态；后续按等深线分七区，
只对浅而平缓的 I、II 区另行拟合。这里应模仿的不是“针对上述问题”六个字，而是旧方案
的具体负担、新方案的下游输入输出关系以及失败区域和修复范围都没有省略。

结果交付同样按阅读任务分流。路线过密时，正文图只说明位置，精确坐标移到附件，
主文继续按区给条数、长度与面积，并汇总 94 条、$434662\,\mathrm{m}$、有效覆盖
$99.956\%$ 和漏测 $0.04\%$。多起始线只支持已比较候选中的代表；连续航线和返航仍留给
Hamilton/Euler 路线输入输出关系。写“取最优方案”以前先说明候选范围，正是人类论文避免事后
夸大结论的一种朴素写法。

B159 与 B195 把抽检题里的选择过程写得很有层次。B159 先问企业更怕把次品放行，还是
更怕把合格批拒收，由这一态度反向确定 $H_0/H_1$，再从功效、样本量和成本进入决策；
B195 面对没有真实抽检样本的条件，先公开模拟器及其假设，待每一步费用、时间和状态转移
闭合后，才压缩重复流程。两篇都不把“采用动态规划”当作起点，而让检验立场、数据条件和
流程复杂度逐级授权后面的模型。

B196 的段落从结构计数起步：节点数、反馈边数和候选动作一旦达到某个规模，穷举负担才
成为换用蚁群或遗传算法的理由。结果解释也没有停在收敛曲线，而是先找反复出现的关键节点，
再说明这些节点对应哪一步真实操作；扰动后则依次回答哪些节点发生变化、总成本怎样变化、
为什么是这一方向。这样的“结构--搜索--热点--执行”比列出三个智能优化算法更有内容。

B060 与 B157 代表另一条机理反演写法。B060 把局部谱线与全曲线反演分工：前者便于确定
候选区间，后者承担最终参数估计；若中间频谱发生变化，还要把差异重新送入厚度模型，不能
停在中间变量。B157 则先说明频率干扰会怎样影响厚度反演，再决定分段修正；当主体厚度
几乎不变而局部残差下降时，作者把贡献限定在局部拟合，不把有限改善扩大成“全面提升”。
这种主动收窄结论强度的语气，是人类技术论文很重要的一种可信感。

